\documentclass[../main.tex]{subfiles}
\graphicspath{{\subfix{../img/}}{\subfix{../architectures images/}}}
\begin{document}

\chapter{Sprint 3 : Contr\^{o}le d'Acc\`{e}s et Int\'{e}gration Mat\'{e}rielle}

\section{Introduction}

Ce sprint est celui qui distingue TrustGym d'une simple application de gestion : la plateforme s'interface avec le mat\'{e}riel physique. Terminaux ZKTeco \cite{zkteco}, cam\'{e}ras Hikvision \cite{hikvision}, badges \acrshort{rfid} --- chaque entr\'{e}e et sortie est captur\'{e}e, v\'{e}rifi\'{e}e et visible en temps r\'{e}el sur le tableau de bord et l'application mobile. Piloter des dispositifs physiques depuis une architecture web exige une couche interm\'{e}diaire d\'{e}di\'{e}e : c'est le r\^{o}le du Machine Handler, connect\'{e} \`{a} l'\acrshort{api} via WebSocket \cite{socketio} et PostgreSQL \texttt{LISTEN/NOTIFY}.

\section{Sprint Backlog}

\begin{longtable}{|p{3cm}|p{2.5cm}|p{4cm}|p{5cm}|}
\hline
\textbf{Epic} & \textbf{Nature} & \textbf{User Story} & \textbf{Description} \\
\hline
\endfirsthead
\hline
\textbf{Epic} & \textbf{Nature} & \textbf{User Story} & \textbf{Description} \\
\hline
\endhead
\hline
\endfoot

\multirow{3}{3cm}{Gestion des Portes} & Fonctionnalit\'{e} & En tant qu'administrateur, je veux configurer une porte connect\'{e}e & Impl\'{e}menter le \acrshort{crud} des portes avec IP, port, type de \acrshort{sdk} (PULL, PUSH-ATTENDANCE, PUSH-SECURITY, HIKVISION, STANDALONE) et num\'{e}ro de s\'{e}rie. \\
\cline{2-4}
& Fonctionnalit\'{e} & En tant qu'administrateur, je veux g\'{e}rer les cr\'{e}neaux horaires & Configurer les plages horaires d'ouverture et les restrictions d'acc\`{e}s par cr\'{e}neau. \\
\cline{2-4}
& Fonctionnalit\'{e} & En tant qu'administrateur, je veux g\'{e}rer les cr\'{e}neaux horaires des partenaires & Permettre aux partenaires de d\'{e}finir leurs propres horaires d'acc\`{e}s. \\
\hline

\multirow{3}{3cm}{\'{E}v\'{e}nements d'Acc\`{e}s} & Fonctionnalit\'{e} & En tant qu'administrateur, je veux voir l'historique d'acc\`{e}s & Enregistrer et afficher les \'{e}v\'{e}nements d'acc\`{e}s avec date, heure, cardId, statut d'autorisation et type (ENTRY/EXIT). \\
\cline{2-4}
& Fonctionnalit\'{e} & En tant que caissier, je veux enregistrer manuellement une entr\'{e}e/sortie & Permettre l'enregistrement manuel d'\'{e}v\'{e}nements d'acc\`{e}s depuis le tableau de bord. \\
\cline{2-4}
& Fonctionnalit\'{e} & En tant qu'administrateur, je veux filtrer l'historique d'acc\`{e}s & Ajouter des filtres avanc\'{e}s par date, membre, porte, type et statut d'autorisation. \\
\hline

\multirow{4}{3cm}{Machine Handler} & T\^{a}che & Mettre en place le service Python & D\'{e}velopper le service Python (port 8000) pour la communication directe avec les dispositifs ZKTeco et Hikvision. \\
\cline{2-4}
& T\^{a}che & Mettre en place le service NestJS relais & D\'{e}velopper le service NestJS (port 8080) pour relayer les \'{e}v\'{e}nements vers \texttt{trustgym-api}. \\
\cline{2-4}
& T\^{a}che & Configurer ngrok et PM2 & Mettre en place le tunnel ngrok \cite{ngrok} pour l'exposition externe et PM2 \cite{pm2} pour la gestion des processus. \\
\cline{2-4}
& T\^{a}che & Cr\'{e}er le lanceur Windows & D\'{e}velopper \texttt{start-trustgym.bat} pour le d\'{e}marrage automatis\'{e} de l'ensemble des services. \\
\hline

\multirow{3}{3cm}{Temps R\'{e}el} & Fonctionnalit\'{e} & En tant qu'administrateur, je veux voir l'occupation en temps r\'{e}el & Calculer l'occupation en comptant les membres dont le dernier \'{e}v\'{e}nement est ENTRY dans une fen\^{e}tre configurable (d\'{e}faut : 1 heure). \\
\cline{2-4}
& T\^{a}che & Impl\'{e}menter le WebSocket & Mettre en place Socket.IO sur le namespace \texttt{/attendance} avec l'\'{e}v\'{e}nement \texttt{occupancy:updated}. \\
\cline{2-4}
& T\^{a}che & Impl\'{e}menter PostgreSQL LISTEN/NOTIFY & Configurer \texttt{PgListenerService} pour la propagation en temps r\'{e}el des \'{e}v\'{e}nements d'acc\`{e}s. \\
\hline

\multirow{3}{3cm}{Commandes et Bundles} & Fonctionnalit\'{e} & En tant qu'administrateur, je veux envoyer des commandes aux dispositifs & Impl\'{e}menter la file de commandes (\texttt{DeviceCommand}) pour piloter les dispositifs ZKTeco/Hikvision. \\
\cline{2-4}
& Fonctionnalit\'{e} & En tant qu'administrateur, je veux suivre l'\'{e}tat des handlers & Cr\'{e}er le tableau de bord des bundles (\texttt{BundleDashboard}) avec statut et flux de configuration. \\
\cline{2-4}
& Fonctionnalit\'{e} & En tant qu'administrateur, je veux mettre \`{a} jour les handlers & Impl\'{e}menter le module de mises \`{a} jour des bundles (\texttt{BundleUpdate}). \\
\hline

\caption{Sprint Backlog -- Sprint 3}
\label{tab:sprint3-backlog}
\end{longtable}

\section{Analyse}

\subsection{Diagramme de cas d'utilisation}

La figure \ref{fig:sprint3-usecase} r\'{e}sume les actions disponibles : l'administrateur configure le mat\'{e}riel et consulte les acc\`{e}s, le caissier peut les enregistrer manuellement, et le membre voit l'occupation en direct.

\begin{figure}[H]
\centering
\begin{tikzpicture}[scale=0.85, transform shape]

% System box
\node[systembox, minimum width=10cm, minimum height=12cm] (sys) at (5.5, -5) {};
\node[uctitle] at (sys.north) {Syst\`{e}me TrustGym -- Contr\^{o}le d'Acc\`{e}s};

% Actors (all on the left)
\stickfigure{(-2.5, -1.5)}{Admin}
\stickfigure{(-2.5, -5)}{Caissier}
\stickfigure{(-2.5, -9)}{Membre}

% Use cases
\node[usecase] (uc1) at (5.5, 0) {Configurer une porte};
\node[usecase] (uc2) at (5.5, -1.8) {G\'{e}rer les cr\'{e}neaux horaires};
\node[usecase] (uc3) at (5.5, -3.6) {Consulter l'historique d'acc\`{e}s};
\node[usecase] (uc4) at (5.5, -5.4) {Enregistrer manuellement un acc\`{e}s};
\node[usecase] (uc5) at (5.5, -7.2) {Consulter l'occupation en temps r\'{e}el};
\node[usecase] (uc6) at (5.5, -9) {Envoyer une commande au dispositif};
\node[usecase] (uc7) at (5.5, -10.8) {Suivre l'\'{e}tat des handlers};

% Admin links
\draw (-1, -1.5) -- (uc1.west);
\draw (-1, -1.5) -- (uc2.west);
\draw (-1, -1.5) -- (uc3.west);
\draw (-1, -1.5) -- (uc5.west);
\draw (-1, -1.5) -- (uc6.west);
\draw (-1, -1.5) -- (uc7.west);

% Caissier links
\draw (-1, -5) -- (uc3.west);
\draw (-1, -5) -- (uc4.west);

% Membre links
\draw (-1, -9) -- (uc5.west);

\end{tikzpicture}
\caption{Diagramme de cas d'utilisation -- Sprint 3}
\label{fig:sprint3-usecase}
\end{figure}

\subsection{Cas d'utilisation cibl\'{e}s}

Format d\'{e}fini en section~\ref{sec:uc_format}. Ce sprint couvre : \textbf{Configurer une porte} (Admin --- IP, port, type de \acrshort{sdk}, test de connectivit\'{e}), \textbf{Consulter l'occupation en temps r\'{e}el} (Admin/Membre, \texttt{occupancy:updated} via WebSocket), et \textbf{Enregistrer manuellement un acc\`{e}s} (Caissier/Admin, \texttt{isAuthorized=false} si le membre est bloqu\'{e}).

\section{R\'{e}alisation}

\subsection{Backend}

\subsubsection{Modules m\'{e}tier}

\begin{longtable}{|>{\footnotesize}p{4.5cm}|>{\footnotesize}p{3.6cm}|p{5.6cm}|}
\hline
\textbf{Module} & \textbf{Entit\'{e}s cl\'{e}s} & \textbf{R\^{o}le} \\
\hline
\endfirsthead
\hline
\endfoot
\texttt{AccessModule} & \texttt{Access} & \'{E}v\'{e}nements ENTRY/EXIT, index\'{e}s sur \texttt{doorId}, \texttt{dateTime}, \texttt{accessType}, \texttt{memberId} ; statut \texttt{isAuthorized}. \\
\hline
\texttt{DoorModule} & \texttt{Door} & Dispositifs : IP, port, num\'{e}ro de s\'{e}rie ; types de \acrshort{sdk} : PULL, PUSH-ATTENDANCE, PUSH-SECURITY, HIKVISION, STANDALONE. \\
\hline
\texttt{ZkecoModule} & --- & Pont entre l'\acrshort{api} et les dispositifs ZKTeco (ouverture, synchronisation des cartes). \\
\hline
\texttt{TimeSlotModule}, \texttt{PartnerTimeSlotModule} & \texttt{TimeSlotHours}, \texttt{PartnerTimeSlot} & Plages horaires d'ouverture par jour, restrictions d'acc\`{e}s. \\
\hline
\texttt{WebSocketModule} (\texttt{AccessGateway}) & --- & Socket.IO sur \texttt{/attendance} : \texttt{authenticate}/\texttt{auth:success}, \texttt{attendance:recorded}, \texttt{occupancy:updated}. Salles \texttt{gym:<gymId>}, aliment\'{e}es par un trigger PostgreSQL \texttt{NOTIFY} consomm\'{e} par \texttt{PgListenerService}. \\
\hline
\caption{Modules backend du Sprint 3}
\end{longtable}

\subsubsection{Calcul de l'Occupation en Direct}

L'occupation en temps r\'{e}el est calcul\'{e}e en comptant les membres dont le dernier \'{e}v\'{e}nement d'acc\`{e}s \`{a} une porte de la salle est de type \texttt{ENTRY}, dans une fen\^{e}tre temporelle configurable (par d\'{e}faut 1 heure). Ce m\'{e}canisme \'{e}vite les faux positifs li\'{e}s aux membres qui auraient quitt\'{e} la salle sans badger en sortie. Chaque nouvel \'{e}v\'{e}nement d\'{e}clenche l'\'{e}v\'{e}nement WebSocket \texttt{occupancy:updated}.

\subsubsection{Algorithme d'Autorisation d'Acc\`{e}s}

L'autorisation d'acc\`{e}s suit un algorithme en six \'{e}tapes ex\'{e}cut\'{e} \`{a} chaque badge \acrshort{rfid} :

\begin{enumerate}
  \item \textbf{Identification} : recherche de la personne par \texttt{cardId} dans l'ordre : Member $\rightarrow$ Coach $\rightarrow$ User.
  \item \textbf{V\'{e}rification du blocage} : si le membre est bloqu\'{e} (\texttt{isBlocked = true}), l'acc\`{e}s est refus\'{e}.
  \item \textbf{V\'{e}rification de l'abonnement actif} : recherche d'un \texttt{MemberGymSubscription} valide (plage de dates ou sessions restantes $> 0$).
  \item \textbf{V\'{e}rification de la porte} : contr\^{o}le via \texttt{GymSubscriptionDoor} que l'abonnement couvre la porte utilis\'{e}e.
  \item \textbf{V\'{e}rification du cr\'{e}neau horaire} : contr\^{o}le via \texttt{SubscriptionTimeSlot} que l'heure actuelle est dans la plage autoris\'{e}e.
  \item \begin{sloppypar}\textbf{D\'{e}cr\'{e}mentation} : si l'abonnement est de type \texttt{PER\_SESSION}, le compteur \texttt{numberOfSessions} est d\'{e}cr\'{e}ment\'{e}.\end{sloppypar}
\end{enumerate}

\begin{sloppypar}
L'enregistrement \texttt{Access} est cr\'{e}\'{e} avec le statut \texttt{isAuthorized} correspondant, puis les \'{e}v\'{e}nements WebSocket sont \'{e}mis.
\end{sloppypar}

\subsubsection{Verrouillage des Dispositifs (DeviceLockService)}

Le \texttt{DeviceLockService} impl\'{e}mente un mutex par dispositif, cl\'{e} sur le couple IP:port. Ce m\'{e}canisme est critique car les \acrshort{sdk}s ZKTeco n'autorisent pas les acc\`{e}s concurrents sur le m\^{e}me port COM ou la m\^{e}me connexion DLL. Sans ce verrouillage, deux requ\^{e}tes simultan\'{e}es vers le m\^{e}me dispositif provoqueraient un crash du service. Le mutex est acquis avant chaque op\'{e}ration mat\'{e}rielle et lib\'{e}r\'{e} automatiquement apr\`{e}s compl\'{e}tion, garantissant l'int\'{e}grit\'{e} des communications.

\subsubsection{Protocole ADMS (Push-Mode EFace10)}

Les dispositifs ZKTeco EFace10 fonctionnent en mode PUSH via le protocole ADMS (\textit{Automatic Data Master Server}). Le Machine Handler expose les endpoints suivants pour ces dispositifs :

\begin{itemize}
  \item \begin{sloppypar}\texttt{POST /iclock/registry} : enregistrement initial du dispositif aupr\`{e}s du serveur.\end{sloppypar}
  \item \begin{sloppypar}\texttt{GET /iclock/push} : configuration des options de polling du dispositif.\end{sloppypar}
  \item \begin{sloppypar}\texttt{GET /iclock/getrequest} : le dispositif interroge le serveur pour r\'{e}cup\'{e}rer les commandes en attente (synchronisation de cartes, red\'{e}marrage).\end{sloppypar}
  \item \begin{sloppypar}\texttt{POST /iclock/cdata} : envoi des donn\'{e}es de pointage depuis le dispositif.\end{sloppypar}
  \item \begin{sloppypar}\texttt{POST /iclock/devicecmd} : acquittement des commandes ex\'{e}cut\'{e}es.\end{sloppypar}
\end{itemize}

\begin{sloppypar}
L'entit\'{e} \texttt{DeviceCommand} g\`{e}re la file de commandes avec un suivi de statut : \texttt{queued} $\rightarrow$ \texttt{sent} $\rightarrow$ \texttt{acked} ou \texttt{failed}.
\end{sloppypar}

\subsubsection{Modules Commandes et Bundles}

\begin{sloppypar}
Le \texttt{DeviceCommandModule} g\`{e}re une file de commandes destin\'{e}es aux dispositifs physiques (synchronisation de cartes, ouverture \`{a} distance, red\'{e}marrage). Les modules \texttt{BundleDashboardModule}, \texttt{BundleSetupModule} et \texttt{BundleUpdateModule} permettent respectivement le suivi de l'\'{e}tat des handlers, leur configuration initiale et leur mise \`{a} jour \`{a} distance.
\end{sloppypar}

\subsection{Machine Handler}

\subsubsection{Architecture}

Communiquer directement avec du mat\'{e}riel ZKTeco depuis NestJS n'est pas possible --- les \acrshort{sdk}s sont en Python (32 bits) et exigent des connexions DLL s\'{e}rielles. D'\`{o} une architecture en deux couches :

\begin{itemize}
  \item \textbf{Service Python} (port 8000) : communique directement avec les dispositifs physiques ZKTeco et Hikvision via leurs \acrshort{sdk}s respectifs. Ce service g\`{e}re la connexion bas niveau, la lecture des \'{e}v\'{e}nements et l'envoi de commandes.
  \item \textbf{Service NestJS} (port 8080) : sert de relais entre le service Python et l'\acrshort{api} principale \texttt{trustgym-api}. Il normalise les donn\'{e}es, g\`{e}re les retentatives en cas d'\'{e}chec et maintient un journal des \'{e}v\'{e}nements.
\end{itemize}

La figure \ref{fig:arch_machines} illustre l'architecture compl\`{e}te du Machine Handler et le flux de donn\'{e}es entre les diff\'{e}rentes couches.

\begin{figure}[H]
  \centering
  \includegraphics[width=\textwidth]{ARCHITECTURE MACHINES.png}
  \caption{Architecture du Machine Handler -- Communication mat\'{e}rielle}
  \label{fig:arch_machines}
\end{figure}

\subsubsection{Infrastructure}

\begin{itemize}
  \item \textbf{ngrok} \cite{ngrok} : cr\'{e}e un tunnel s\'{e}curis\'{e} \acrshort{https} exposant le service NestJS (port 8080) sur Internet via un domaine r\'{e}serv\'{e} par salle (ex. \texttt{samurai.trustgym.net}), permettant aux dispositifs install\'{e}s en mode PUSH d'envoyer leurs \'{e}v\'{e}nements m\^{e}me derri\`{e}re un \acrshort{nat}.
  \item \textbf{PM2} \cite{pm2} : gestionnaire de processus utilis\'{e} lors de la phase de d\'{e}veloppement pour orchestrer et superviser les services Python et NestJS avec red\'{e}marrage automatique.
  \item \textbf{NSSM} \cite{nssm} (\textit{Non-Sucking Service Manager}) : solution retenue pour la production, permettant d'enregistrer les processus Python et NestJS comme de v\'{e}ritables \textbf{Services Windows} persistants, ind\'{e}pendants de toute session utilisateur et d\'{e}marrant automatiquement avec le syst\`{e}me d'exploitation.
\end{itemize}

%──────────────────────────────────────────────
\subsubsection{Installateur Windows (Inno Setup)}
%──────────────────────────────────────────────

\begin{sloppypar}
Pour simplifier le d\'{e}ploiement du Machine Handler sur les \acrshort{pc} des salles partenaires, un installateur Windows autonome a \'{e}t\'{e} cr\'{e}\'{e} \`{a} l'aide d'\textbf{Inno Setup 6} \cite{innosetup}. Le script \texttt{TrustGymSetup.iss} g\'{e}n\`{e}re un ex\'{e}cutable \texttt{TrustGymSetup-1.0.0.exe} qui empaquette l'int\'{e}gralit\'{e} du bundle -- runtimes inclus -- dans un seul fichier distribuable par cl\'{e} \acrshort{usb} ou t\'{e}l\'{e}chargement direct.
\end{sloppypar}

\paragraph{Deux modes de configuration.}
L'assistant propose deux modes d\'{e}terminant la mani\`{e}re dont les param\`{e}tres sont collect\'{e}s :

\begin{itemize}
  \item \textbf{Mode Token (recommand\'{e})} : le Super Admin g\'{e}n\`{e}re un jeton de configuration unique (format \texttt{TG-XXXX-XXXX}) depuis le tableau de bord cloud. L'installateur appelle alors l'endpoint \url{GET /api/bundle-setup/config?token=...} et r\'{e}cup\`{e}re automatiquement tous les param\`{e}tres : type de dispositif, token ngrok, domaine, identifiant de la salle, cl\'{e}s \acrshort{api} et informations de connexion \`{a} la base de donn\'{e}es. Le technicien n'a qu'\`{a} saisir le jeton et cliquer sur \og{} Suivant \fg{}.
  \item \textbf{Mode Manuel} : le technicien saisit manuellement chaque param\`{e}tre \`{a} travers les \'{e}crans successifs de l'assistant : type de dispositif, token et domaine ngrok, identifiant de la salle (\acrshort{uuid}), cl\'{e} \acrshort{api} bridge et environnement de base de donn\'{e}es (production, staging ou custom).
\end{itemize}

\paragraph{Contenu du package.}
\begin{sloppypar}
L'installateur copie l'ensemble des composants dans \texttt{C:\textbackslash{}Program Files (x86)\textbackslash{}TrustGym\textbackslash{}} :
\end{sloppypar}

\begin{itemize}
  \item Le service NestJS compil\'{e} (\texttt{trustgym-zktco-api/dist/})
  \item Le service Python (Hikvision, ZKTeco C3-200, ZKTeco SCR100 ou ZKTeco SC105 selon le profil)
  \item Les runtimes embarqu\'{e}s : Python 3.13 (32 bits), Node.js 20.18 et NSSM 2.24
  \item L'ex\'{e}cutable ngrok et les scripts PowerShell :
        \texttt{unified-install.ps1}, \texttt{unified-uninstall.ps1}, \texttt{unified-watchdog.ps1}
\end{itemize}

\paragraph{G\'{e}n\'{e}ration des fichiers de configuration.}
\begin{sloppypar}
Apr\`{e}s la copie des fichiers, l'installateur g\'{e}n\`{e}re dynamiquement deux fichiers de configuration via un script PowerShell temporaire \'{e}crit dans \texttt{\%TEMP\%} (contournement n\'{e}cessaire car Inno Setup ne peut pas \'{e}crire directement dans \texttt{Program Files}) :
\end{sloppypar}
\begin{itemize}
  \item \texttt{env} : contient le token ngrok, le domaine r\'{e}serv\'{e}, le type de \acrshort{sdk} et les chemins vers les runtimes embarqu\'{e}s.
  \item \texttt{trustgym-zktco-api/.env} : contient les coordonn\'{e}es de la base de donn\'{e}es PostgreSQL, l'identifiant de la salle, les cl\'{e}s \acrshort{api} et la configuration \acrshort{smtp}.
\end{itemize}

\begin{sloppypar}
Une fois la configuration g\'{e}n\'{e}r\'{e}e, l'installateur ex\'{e}cute automatiquement le script \texttt{unified-install.ps1} en mode silencieux, qui prend en charge la cr\'{e}ation des Services Windows et le d\'{e}marrage des processus.
\end{sloppypar}

\begin{figure}[H]
  \centering
  \includegraphics[width=\textwidth, height=0.75\textheight, keepaspectratio]{capture/Assistant-Inno-Setup.png}
  \caption{Assistant Inno Setup -- s\'{e}lection du mode de configuration}
  \label{fig:sprint3-innosetup-mode}
\end{figure}

%──────────────────────────────────────────────
\subsubsection{Services Windows Persistants (NSSM)}
%──────────────────────────────────────────────

Le script \texttt{unified-install.ps1} utilise \textbf{NSSM 2.24} pour enregistrer les processus du Machine Handler comme des Services Windows natifs. Contrairement \`{a} PM2 -- qui n\'{e}cessite une session utilisateur ouverte -- les services NSSM d\'{e}marrent d\`{e}s l'allumage du \acrshort{pc}, sans aucune interaction humaine, et se relancent automatiquement en cas de plantage.

\paragraph{Services cr\'{e}\'{e}s.}
Le tableau~\ref{tab:nssm_services} r\'{e}capitule les quatre services install\'{e}s sur chaque salle partenaire :

\begin{table}[H]
\centering
\small
\begin{tabular}{|p{4.5cm}|p{2cm}|p{2.5cm}|p{4.5cm}|}
\hline
\textbf{Service} & \textbf{Gestionnaire} & \textbf{Port} & \textbf{Comportement} \\
\hline
\texttt{TrustGym-Python-API} & NSSM & 8000 / 8001 / 8002 & D\'{e}marrage automatique, redémarrage $\times$3 (d\'{e}lai 5 s) \\
\hline
\texttt{TrustGym-NestJS-API} & NSSM & 8080 & D\'{e}marrage diff\'{e}r\'{e} (\textit{DelayedAutoStart}), d\'{e}pend de Python API \\
\hline
\texttt{ngrok} & Natif ngrok & -- & Tunnel \acrshort{https} vers le domaine r\'{e}serv\'{e} de la salle \\
\hline
\texttt{TrustGym-Watchdog} & T\^{a}che planifi\'{e}e & -- & V\'{e}rification des endpoints \texttt{/health} toutes les 5 min, red\'{e}marrage automatique en cas d'\'{e}chec \\
\hline
\end{tabular}
\caption{Services Windows install\'{e}s par NSSM sur chaque salle partenaire}
\label{tab:nssm_services}
\end{table}

\paragraph{Configuration NSSM.}
\begin{sloppypar}
Pour chaque service, le script configure les propri\'{e}t\'{e}s suivantes via des appels successifs \`{a} \texttt{nssm.exe set} :
\end{sloppypar}

\begin{itemize}
  \item \textbf{R\'{e}pertoire de travail} (\texttt{AppDirectory}) : r\'{e}pertoire d'installation du service.
  \item \begin{sloppypar}\textbf{Journalisation} (\texttt{AppStdout} / \texttt{AppStderr}) : fichiers de log s\'{e}par\'{e}s dans \texttt{C:\textbackslash{}ProgramData\textbackslash{}TrustGym\textbackslash{}logs\textbackslash{}}, avec rotation automatique \`{a} 10~Mo (\texttt{AppRotateBytes}).\end{sloppypar}
  \item \begin{sloppypar}\textbf{Variables d'environnement} (\texttt{AppEnvironmentExtra}) : injection de \texttt{NODE\_ENV=production}, \texttt{HOST=127.0.0.1} et \texttt{PORT=8080} pour le service NestJS.\end{sloppypar}
  \item \textbf{Politique de reprise} (\texttt{sc failure}) : red\'{e}marrage automatique trois fois cons\'{e}cutives avec un d\'{e}lai de 5 secondes entre chaque tentative, puis r\'{e}initialisation du compteur apr\`{e}s 24 heures.
  \item \begin{sloppypar}\textbf{D\'{e}marrage diff\'{e}r\'{e}} : la cl\'{e} de registre \texttt{DelayedAutoStart} est activ\'{e}e pour \texttt{TrustGym-NestJS-API} afin d'\'{e}viter les conflits au d\'{e}marrage du syst\`{e}me (\textit{boot storm}).\end{sloppypar}
\end{itemize}

\paragraph{Watchdog.}
\begin{sloppypar}
Une t\^{a}che planifi\'{e}e Windows (\textit{Scheduled Task}) nomm\'{e}e \texttt{TrustGym-Watchdog} est cr\'{e}\'{e}e et ex\'{e}cut\'{e}e toutes les cinq minutes. Elle interroge les endpoints \texttt{/health} de chaque service et red\'{e}marre ceux qui ne r\'{e}pondent pas. Ce m\'{e}canisme compl\`{e}te la politique de red\'{e}marrage NSSM en d\'{e}tectant les situations de gel (\textit{deadlock}) o\`{u} le processus est actif mais non fonctionnel.
\end{sloppypar}

\subsection{Diagramme de s\'{e}quence}

La figure \ref{fig:sprint3-sequence} trace le chemin d'un badge \acrshort{rfid} jusqu'au tableau de bord : six \'{e}tapes, sous la seconde, de bout en bout.

\begin{figure}[H]
\centering
\begin{tikzpicture}[scale=0.78, transform shape]

% Participants
\node[sdactor] (dev) at (0, 0) {Dispositif};
\node[sdactor] (py) at (3, 0) {Python 8000};
\node[sdactor] (nest) at (6, 0) {NestJS 8080};
\node[sdactor] (api) at (9, 0) {trustgym-api};
\node[sdactor] (pg) at (12, 0) {PostgreSQL};
\node[sdactor] (ws) at (15, 0) {Dashboard};

% Lifelines
\draw[dashed] (0, -0.35) -- (0, -13);
\draw[dashed] (3, -0.35) -- (3, -13);
\draw[dashed] (6, -0.35) -- (6, -13);
\draw[dashed] (9, -0.35) -- (9, -13);
\draw[dashed] (12, -0.35) -- (12, -13);
\draw[dashed] (15, -0.35) -- (15, -13);

% Messages
\draw[sdmessage] (0, -1) -- node[above, font=\scriptsize] {Badge RFID} (3, -1);
\draw[sdmessage] (3, -2) -- node[above, font=\scriptsize] {POST /event} (6, -2);
\draw[sdmessage] (6, -3) -- node[above, font=\scriptsize] {POST /api/access} (9, -3);
\draw[sdmessage] (9, -4) -- node[above, font=\scriptsize] {INSERT access} (12, -4);
\draw[sdreturn] (12, -5) -- node[above, font=\scriptsize] {OK} (9, -5);

% LISTEN/NOTIFY flow
\draw[sdmessage] (12, -6.5) -- node[above, font=\scriptsize] {NOTIFY access\_log\_channel} (9, -6.5);
\node[font=\scriptsize\itshape, text=gray] at (10.5, -7.2) {PgListenerService};

\draw[sdmessage] (9, -8) -- node[above, font=\scriptsize] {Calcul occupation} (12, -8);
\draw[sdreturn] (12, -9) -- node[above, font=\scriptsize] {count: 42} (9, -9);

\draw[sdmessage] (9, -10.5) -- node[above, font=\scriptsize, text width=3.5cm, align=center] {emit occupancy:updated} (15, -10.5);
\node[font=\scriptsize\itshape, text=gray] at (12, -11.2) {Socket.IO /attendance};

\draw[sdreturn] (6, -12) -- node[above, font=\scriptsize] {200 OK} (3, -12);

\end{tikzpicture}
\caption{Diagramme de s\'{e}quence -- Flux temps r\'{e}el d'un \'{e}v\'{e}nement d'acc\`{e}s}
\label{fig:sprint3-sequence}
\end{figure}

\subsection{Interfaces r\'{e}alis\'{e}es (web et mobile)}

Pages livr\'{e}es : \texttt{historique} (filtres dates/membre/porte/type), \texttt{enregistrer} (acc\`{e}s manuel), \texttt{horaire} (plages d'ouverture). C\^{o}t\'{e} mobile, le hook \texttt{useOccupancyRealtime} combine Socket.IO (\texttt{occupancy:updated}) et un polling de secours toutes les 30~s, mis en pause sur perte de r\'{e}seau pour pr\'{e}server la batterie.

\begin{figure}[H]
  \centering
  \includegraphics[width=\textwidth, height=0.55\textheight, keepaspectratio]{capture/historique-acces.png}
  \caption{Historique des acc\`{e}s avec filtres avanc\'{e}s}
  \label{fig:sprint3-access-history}
\end{figure}

\begin{figure}[H]
  \centering
  \includegraphics[width=\textwidth, height=0.55\textheight, keepaspectratio]{capture/config-portes.png}
  \caption{Configuration des portes et dispositifs de contr\^{o}le d'acc\`{e}s}
  \label{fig:sprint3-doors}
\end{figure}

\begin{figure}[H]
  \centering
  \includegraphics[width=\textwidth, height=0.55\textheight, keepaspectratio]{capture/Toast entry-exit socketio.jpeg}
  \caption{Occupation en direct sur l'\'{e}cran mobile}
  \label{fig:sprint3-mobile-occupancy}
\end{figure}

\section{Conclusion}

Le Machine Handler, l'algorithme d'autorisation en six \'{e}tapes, le tunnel ngrok, les services Windows NSSM --- ce sprint a \'{e}t\'{e} le plus exigeant techniquement. Mais il est aussi celui qui donne \`{a} TrustGym sa valeur ajout\'{e}e r\'{e}elle : une porte qui s'ouvre parce que le syst\`{e}me sait qui est devant. Le Sprint 4 reprend sur cette base pour ajouter les coachs et leurs s\'{e}ances.

\end{document}
